\begin{document}
\section{Supplementary material}
Figure S1: Spatial variograms of the binary component. Each color represents a different scenario. A total of 40 variograms is produced, one per time step. A limit of 1000 km is imposed for the spatial correlation range.

Figure S2: Spatial variograms of the continuous component. Only one variogram is computed per scenario. A limit of 500 km is imposed for the spatial correlation range.

Figure S3: Temporal variograms of the continuous component. Only one variogram is computed per scenario. No temporal correlation limit is imposed.

Figure S4: Full five-fold and hindcasting cross-validations results. 

Figure S5: Results of RelAb \~ PoO(VarSel\_IK) + climate(VarSel\_STRK). This output is used in the main manuscript for the description of the post-glacial expansion of \textit{Abies} spp. 

Figure S6: Separability index calculated under each scenario.

Figure S7: Percentage of missing data under different temporal bins over time. 

Figure S8: Pollen sum per sample. Left: pollen sums per sample not temporally binned; center: pollen sums per sample with a bin size of 500 years; right: pollen sums per sample with a bin size of 200 years. The two horizontal lines indicate the 300 and 1000 pollen sums used as the minimum requirements in this study. 

Figure S9: Sensitivity test of the binary components to different time bin size. This figure show cross validation results for the four scenarios used in this work. Different colors represents 200 (yellow) and 500 (black) time bin size. The interpretation is equivalent of Fig. 4A.

Table S1: Site metadata. Information is given about the sites present in the original dataset (Lang et al. 2023) and in the manually added one. In the latter case, the origin, related publications, and calibration model are specified.

Table S2: Correction factors applied to raw pollen data. For more information refer to Lang et al. (2023; page 153).

Supplementary review: review of \textit{Abies} spp. demographic history in Europe by comparing our result with past literature. 

\section{References}

Lang, G., Ammann, B., Behre, K.-E. and Tinner, W. (2023) Quaternary vegetation dynamics of Europe. Haupt Verlag.

Oriani, F., Mariethoz, G. and Chevalier, M. (2024) Eupollmap: the European atlas of contemporary pollen distribution maps derived from an integrated kriging interpolation approach. Earth System Science Data, 16, 731–742. URL: http://dx.doi.org/10.5194/essd-16-731-2024.
\end{document}